% block: direct-methods-introduction pages:17,18 type:section
\subsection{{Direct Methods}}\label{sec:direct-methods}
{For simplicity of presentation, we only consider the case $A\bm x=\bm b$ where $A\in\mathbb R^{n\times n}$ and $\bm b\in\mathbb R^n$. In Section~\ref{subsec:lu-decomposition}, the LU decomposition is introduced to solve linear systems, and then in Section~\ref{subsec:lu-pivoting} the LU decomposition with pivoting techniques is explained, which is numerically stable and widely used in practice. A way to improve the accuracy of the approximate solution is described in Section~\ref{subsec:iterative-refinement}.}
% block: lu-factorization page:18 type:subsection
\subsubsection{{LU Decomposition}}\label{subsec:lu-decomposition}
{Let $A$ be a square matrix. Then the LU decomposition of $A$ is achieved by decomposing $A$ into the multiplication of a lower triangular matrix $L$ and an upper triangular matrix $U$ of the form}
\begin{equation}\label{eq:lu-factorization}A=LU,\end{equation}
{where}
\[
L=\begin{pmatrix}l_{11}&&\\\vdots&\ddots&\\l_{n1}&\cdots&l_{nn}\end{pmatrix},\quad
U=\begin{pmatrix}u_{11}&\cdots&u_{1n}\\&\ddots&\vdots\\&&u_{nn}\end{pmatrix}.
\]
{The blanks in matrices $L$ and $U$ mean that all the corresponding elements are zero.}
% block: triangular-substitutions page:18 type:ordinary-paragraph
{If the diagonal elements of $L$ or $U$ are one, then the LU decomposition is uniquely determined. The LU decomposition with $l_{ii}=1$ for all $i$ is referred to as Doolittle's method, and the one with $u_{ii}=1$ for all $i$ is referred to as Crout's method. Once we obtain the LU decomposition of $A$, it is easy to obtain the solution $\bm x$ of the linear systems via}
\begin{equation}\label{eq:lu-triangular-solves}A=LU,\quad L\bm y=\bm b,\quad U\bm x=\bm y.\end{equation}
{To be specific, forward substitution:}
\[
y_i=\frac1{l_{ii}}\left(b_i-\sum_{j=1}^{i-1}l_{ij}y_j\right)\quad(i=1,\ldots,n)
\]
{gives $\bm y$, and then back substitution}
\[
x_i=\frac1{u_{ii}}\left(y_i-\sum_{j=i+1}^n u_{ij}x_j\right)\quad(i=n,\ldots,1)
\]
{yields the solution $\bm x$, where $y_1=b_1/l_{11}$ and $x_n=y_n/u_{nn}$.}
% block: lu-existence pages:18,19 type:ordinary-paragraph
{The necessary and sufficient condition for the existence of the LU decomposition is that all the leading principal minors\footnote{Using MATLAB notation, all the leading principal minors of $A$ are defined by $\det(A(1:1,1:1))$, $\det(A(1:2,1:2))$, \ldots, $\det(A(1:n,1:n))$.} are not zero. Below are some examples of cases where the LU decomposition exists.}
\begin{enumerate}
\item {Strictly diagonally dominant matrix:} $|a_{ii}|>\sum_{j\ne i}|a_{ij}|\quad(1\le i\le n)$.
\item {M-matrix: $a_{ij}\le0\ (i\ne j)$ and all the elements of $A^{-1}$ are non-negative.}
\item {Symmetric part of matrix $A\in\mathbb R^{n\times n}$ is positive definite. The term symmetric part means $(A+A^\top)/2$ and the term positive definite means $\bm v^\top A\bm v>0$ for any nonzero vector $\bm v$. For any square matrix $A$, matrix $A$ can be written as the sum of the symmetric part $(A+A^\top)/2$ and the skew symmetric part $(A-A^\top)/2$.}
\end{enumerate}
% block: algorithm-lu page:19 type:ordinary-paragraph
{The algorithm of the LU decomposition (Doolittle's method) is shown in Algorithm~\ref{alg:lu-decomposition}, which produces $L$ and $U$ such that $A=LU$.}
\begin{algorithm}[H]
\caption{{The LU decomposition}}\label{alg:lu-decomposition}
\begin{algorithmic}[1]
\Statex \textbf{Input:} {an $n\times n$ matrix $A$}
\Statex \textbf{Output:} {a unit lower triangular matrix $L$ and an upper triangular matrix $U$}
\For{$i=1,\ldots,n$}
\State $a_{ji}=a_{ji}/a_{ii}\quad(j=i+1,\ldots,n)$
\For{$j=i+1,\ldots,n$}
\State $a_{jk}=a_{jk}-a_{ji}\times a_{ik}\quad(k=i+1,\ldots,n)$
\EndFor
\EndFor
\State $l_{ii}=1\quad(i=1,\ldots,n)$
\State $l_{ij}=a_{ij}\quad(i>j)$
\State $u_{ij}=a_{ij}\quad(i\le j)$
\end{algorithmic}
\end{algorithm}
% block: pivoting-example pages:19,20 type:subsection
\subsubsection{{LU Decomposition with Pivoting}}\label{subsec:lu-pivoting}
{The LU decomposition may not exist for some $A$. For example, consider}
\[
A=\begin{pmatrix}0&1\\2&3\end{pmatrix}.
\]
{Then it is easy to see that the LU decomposition (Algorithm~\ref{alg:lu-decomposition}) fails due to zero division by $a_{11}=0$. On the other hand, if the first row and the second row are swapped, i.e.,}
\[
PA=\begin{pmatrix}0&1\\1&0\end{pmatrix}\begin{pmatrix}0&1\\2&3\end{pmatrix}
=\begin{pmatrix}2&3\\0&1\end{pmatrix},
\]
{then the LU decomposition does not fail, and we have the LU decomposition of the form}
\[PA=LU,\]
{where}
\[
L=\begin{pmatrix}1&0\\0&1\end{pmatrix},\quad U=\begin{pmatrix}2&3\\0&1\end{pmatrix}.
\]
{This technique is called pivoting.}
% block: partial-pivoting page:20 type:ordinary-paragraph
{The LU decomposition with the partial pivoting is shown in Algorithm~\ref{alg:lu-partial-pivoting}. Here, the term partial pivoting corresponds to lines 2--6 of Algorithm~\ref{alg:lu-partial-pivoting}.}
\begin{algorithm}[H]
\caption{{The LU decomposition with partial pivoting}}\label{alg:lu-partial-pivoting}
\begin{algorithmic}[1]
\Statex \textbf{Input:} {an $n\times n$ matrix $A$ and the $n\times n$ identity matrix $P(=I)$}
\Statex \textbf{Output:} {a unit lower triangular matrix $L$ and an upper triangular matrix $U$}
\Statex \textbf{Output:} {a permutation matrix $P$}
\For{$i=1,\ldots,n$}
\State {Find the maximum element $a_{p_i i}$ from $\{|a_{ii}|,\ldots,|a_{ni}|\}$.}
\State {Swap $a_{ij}$ and $a_{p_i j}$ for $j=1,\ldots,n$.}
\If{$i\ne n$}
\State {Swap $i$th row and $p_i$th row of matrix $P$.}
\EndIf
\State $a_{ji}=a_{ji}/a_{ii}\quad(j=i+1,\ldots,n)$
\For{$j=i+1,\ldots,n$}
\State $a_{jk}=a_{jk}-a_{ji}\times a_{ik}\quad(k=i+1,\ldots,n)$
\EndFor
\EndFor
\State $l_{ii}=1\quad(i=1,\ldots,n)$
\State $l_{ij}=a_{ij}\quad(i>j)$
\State $u_{ij}=a_{ij}\quad(i\le j)$
\end{algorithmic}
\end{algorithm}
{Algorithm~\ref{alg:lu-partial-pivoting} produces $L,U$, and permutation matrix $P$ such that}
\[PA=LU.\]
{Unlike the LU decomposition, the LU decomposition with the partial pivoting never suffers from breakdown (zero-division).}
% block: complete-pivoting pages:20,21 type:ordinary-paragraph
{Next, the LU decomposition with the complete pivoting is shown in Algorithm~\ref{alg:lu-complete-pivoting}. The term complete pivoting corresponds to lines 2--8 of Algorithm~\ref{alg:lu-complete-pivoting}.}
{Notice that at the $i$th step, the partial pivoting finds the maximum value from $\{|a_{k,i}|:k\in\{i,\ldots,n\}\}$ and the complete pivoting finds the maximum value from $\{|a_{k,l}|:k,l\in\{i,\ldots,n\}\}$.}
{Algorithm~\ref{alg:lu-complete-pivoting} produces decomposed matrices $L$ and $U$ and permutation matrices $P$ and $Q$ such that}
\[PAQ=LU.\]
{The LU decomposition with partial pivoting is usually used, and if the accuracy of the decomposition is not satisfactory, then the complete pivoting may be used.}
\begin{algorithm}[H]
\caption{{The LU decomposition with complete pivoting}}\label{alg:lu-complete-pivoting}
\begin{algorithmic}[1]
\Statex \textbf{Input:} {an $n\times n$ matrix $A$ and the $n\times n$ identity matrices $P(=I),Q(=I)$}
\Statex \textbf{Output:} {a unit lower triangular matrix $L$ and an upper triangular matrix $U$}
\Statex \textbf{Output:} {permutation matrices $P,Q$}
\For{$i=1,\ldots,n$}
\State {Find the maximum element $a_{p_iq_i}$ from $\{|a_{kl}|:i\le k\le n,i\le l\le n\}$.}
\State {Swap $a_{ij}$ and $a_{p_i j}$ for $j=1,\ldots,n$.}
\State {Swap $a_{ji}$ and $a_{jq_i}$ for $j=1,\ldots,n$.}
\If{$i\ne n$}
\State {Swap $i$th row and $p_i$th row of matrix $P$.}
\State {Swap $i$th column and $q_i$th row of matrix $P$.}
\EndIf
\State $a_{ji}=a_{ji}/a_{ii}\quad(j=i+1,\ldots,n)$
\For{$j=i+1,\ldots,n$}
\State $a_{jk}=a_{jk}-a_{ji}\times a_{ik}\quad(k=i+1,\ldots,n)$
\EndFor
\EndFor
\State $l_{ii}=1\quad(i=1,\ldots,n)$
\State $l_{ij}=a_{ij}\quad(i>j)$
\State $u_{ij}=a_{ij}\quad(i\le j)$
\end{algorithmic}
\end{algorithm}
% block: multiple-right-hand-sides page:21 type:ordinary-paragraph
{If we need to solve the set of linear systems of the form}
\[A\bm x_i=\bm b_i,\quad i=1,2,\ldots,m,\]
{then the LU decomposition is very attractive because once we obtain $L$ and $U$, then all we have to do is to conduct forward and back substitutions, see the explanation after \eqref{eq:lu-triangular-solves}.}
% block: iterative-refinement pages:21,22 type:subsection
\subsubsection{{Iterative Refinement}}\label{subsec:iterative-refinement}
{In this subsection, we consider the case where the approximate solution $\tilde{\bm x}$ of the linear systems \eqref{eq:linear-system} by the LU decomposition is not accurate enough.}
{In such a case, we can obtain the exact solution $\bm x$ if we know the correction vector $\Delta\bm x$ such that $\bm x=\tilde{\bm x}+\Delta\bm x$. It is easy to see that the correction vector $\Delta\bm x$ corresponds to the solution of the following linear systems:}
\begin{equation}\label{eq:iterative-refinement-correction}A\Delta\bm x=\bm r,\end{equation}
{where $\bm r:=\bm b-A\tilde{\bm x}$. After solving \eqref{eq:iterative-refinement-correction} to obtain an approximate solution $\Delta\tilde{\bm x}$ of \eqref{eq:iterative-refinement-correction}, the corrected approximate solution $\tilde{\bm x}+\Delta\tilde{\bm x}$ is expected to be better in accuracy than $\tilde{\bm x}$. This process can be repeated until the required accuracy is obtained, which is called iterative refinement. The algorithm of the iterative refinement is shown in Algorithm~\ref{alg:iterative-refinement}.}
\begin{algorithm}[H]
\caption{{Iterative refinement for $A\bm x=\bm b$}}\label{alg:iterative-refinement}
\begin{algorithmic}[1]
\Statex \textbf{Input:} $A,\bm b$, {and an approximate solution $\tilde{\bm x}$ for $A\bm x=\bm b$}
\Statex \textbf{Output:} {a refined approximate solution $\tilde{\bm x}$}
\State {Compute $\bm r=\bm b-A\tilde{\bm x}$ using double precision arithmetic.}
\State {Solve $A\Delta\bm x=\bm r$ to obtain the approximate solution $\Delta\tilde{\bm x}$.}
\State {Update $\tilde{\bm x}=\tilde{\bm x}+\Delta\tilde{\bm x}$.}
\State {Repeat steps 1--3 until required accuracy is obtained.}
\end{algorithmic}
\end{algorithm}
% block: refinement-cost-and-newton page:22 type:ordinary-paragraph
{If the LU decomposition is used for $A\bm x=\bm b$, then it is easy to obtain the approximate solution $\Delta\tilde{\bm x}$ because one can reuse the decomposed factors $L$ and $U$, and thus only forward and back substitutions are required.}
{Notice that double precision arithmetic at step 1 in Algorithm~\ref{alg:iterative-refinement} means that if we use single precision arithmetic, then compute $\bm r=\bm b-A\tilde{\bm x}$ using double precision arithmetic, and if we use double precision arithmetic, then compute it using quadruple precision arithmetic.}
{Algorithm~\ref{alg:iterative-refinement} can be interpreted as Newton's method by letting $\bm f(\bm x):=\bm b-A\bm x$ and applying Newton's method to $\bm f(\bm x)=\bm0$.}
